\documentclass[11pt]{article}

\usepackage[a4paper,margin=1in]{geometry}
\usepackage{booktabs}
\usepackage{microtype}
\usepackage{amsmath}
\usepackage{amssymb}
\usepackage{hyperref}
\usepackage{xcolor}

\hypersetup{
  colorlinks=true,
  linkcolor=blue!50!black,
  citecolor=blue!50!black,
  urlcolor=blue!50!black
}

\newcommand{\Ftwo}{\mathbb{F}_2}
\newcommand{\Fp}{\mathbb{F}_p}
\newcommand{\us}{\ensuremath{\mu\mathrm{s}}}
\newcommand{\slred}{\texttt{same-level}}
\newcommand{\ratioiqr}[3]{\shortstack[c]{#1\\[-0.25ex]{\tiny #2--#3}}}

\title{Experimental Notes on Morse Persistence}
\author{Draft note for discussion}
\date{June 16, 2026}

\begin{document}
\maketitle

\begin{abstract}
This note summarizes the current experimental status of the Morse persistence prototype.
The main message is now twofold.  First, the compact C++ backend is already useful for
comparing Morse-sequence strategies, and on selected Roadmap Rips averages it is faster
than the GUDHI baseline used here.  Second, the native \texttt{Gudhi::Simplex\_tree}
integration path has become competitive on structured plateau grids: in the current
30-repeat Release benchmark the direct view beats or essentially matches GUDHI on the
medium grid cases, while GUDHI remains faster on the flag/Rips-like case and on the
largest grid.  The reducer kernel is consistently faster than GUDHI persistence once the
Morse input has been built, but view and Morse-frame construction still dominate the
integrated path.  Across the strategy experiments, critical-cell ratio alone is not a
reliable performance target; annotation volume, mutation work, and construction overhead
explain the observed timings more consistently.  For strategy selection,
\texttt{adaptive} remains the conservative control, while the structured-grid extension
\texttt{adaptive-structured} has become a serious candidate for the next validation round
rather than merely a curiosity.
\end{abstract}

\section{Executive summary}

The goal is not yet to make universal performance claims, but to identify which
implementation choices and algorithmic directions are now supported by measurements.  The
experiments were designed to answer five practical questions:
\begin{enumerate}
  \item Can a Morse-sequence based pipeline reproduce the same barcodes as a standard
        persistence reduction?
  \item Is the number of critical simplices a reliable predictor of performance?
  \item How does the current prototype compare with GUDHI/CAM-style persistence on
        classical Rips benchmarks?
  \item What is the cost of starting directly from a \texttt{Gudhi::Simplex\_tree}, as a
        possible path toward GUDHI integration?
  \item Should we pursue a coreference/cohomological direction, or keep using reference
        annotations after constructing a same-level reduction sequence?
\end{enumerate}

The main conclusion is that the current prototype is already useful for comparing
strategies, while the integration result should be stated carefully.  The compact backend
is faster than GUDHI on the selected Roadmap averages.  The native GUDHI-view path has
become a credible integration route because it avoids the large compact-import cost.  In
the current large Release benchmark it is faster than GUDHI on the smaller structured
grid, at parity on the medium grid, and slower on the flag/Rips-like case and largest
grid.  The technical lesson is unchanged but sharper: critical-cell ratio alone is not
the right optimization target.  Annotation volume, reducer mutation cost, and
frame-construction overhead are more predictive of the observed running time.

\paragraph{Current takeaways.}
\begin{itemize}
  \item Correctness is checked systematically against standard reduction, with selected
        GUDHI validation and optional Perseus barcode export.
  \item On selected Roadmap Rips objects, \slred and
        \texttt{f-max} are the most competitive current strategies in the current
        portfolio artifact, even though
        \slred leaves many more cells critical.
  \item The direct \texttt{Gudhi::Simplex\_tree} view is now a credible integration path
        because it is much faster than importing the tree into the compact complex.
        In the new large 30-repeat native benchmark, the median
        \texttt{GUDHI/Direct} ratio is above $1$ on the $64^2$ plateau grid,
        approximately $1$ on the $96^2$ plateau grid, and below $1$ on the
        flag/Rips-like case and the $128^2$ plateau grid.
  \item On synthetic size-$48$ examples, \texttt{f-max} is the fastest of the four core
        same-level strategies in the representative run, so it must remain in the measured
        portfolio.
  \item The coreference recurrence is correct but too expensive on dense Rips complexes;
        reference reduction remains the useful direction for now.
  \item The selector experiments are now generated from CSV artifacts.  The conservative
        \texttt{adaptive} selector remains the reported control, while
        \texttt{adaptive-structured} eliminates the current image-grid gap in the aggregate
        audit and should be the next selector variant to validate.
  \item The main timing tables below are still over $\Ftwo$ unless stated otherwise, but
        the C++ core and Python interface now also support prime fields $\Fp$; validation
        over $\mathbb{F}_3$ and $\mathbb{F}_5$ shows no barcode mismatches on the tested
        synthetic cases.
\end{itemize}

\tableofcontents

\section{Experimental setup}

\subsection{Implementation used in the experiments}

The implementation is a C++ core with a Python interface.  The Python layer mirrors the C++
objects, but when the nanobind extension is available the finalized complex and Morse
sequence carry stateful C++ objects.  All timings reported here use the C++ backend.

\paragraph{Filtered complex representation.}
The prototype stores an explicit filtered abstract simplicial complex.  Each simplex is
assigned a compact integer identifier.  The finalized representation stores:
\begin{itemize}
  \item the sorted vertex list of each simplex;
  \item the dimension and scalar filtration value;
  \item the discrete filtration level, i.e. the rank of the filtration value among all
        distinct values;
  \item the codimension-one boundary as simplex identifiers;
  \item the coboundary as simplex identifiers;
  \item buckets of simplices by filtration level;
  \item a filtration order compatible with boundary-before-coface reduction.
\end{itemize}
This is not a simplex tree.  At this stage it is deliberately simple: after finalization,
most computations are vector-based and use simplex identifiers rather than vertex tuples.
This makes it convenient to implement Morse steps, reference maps, coreference maps, and
instrumented reductions.  Most performance experiments in this note use $\Ftwo$, because
that is the coefficient field used for the current GUDHI-facing comparison, but the same
core data structure is also used by the prime-field reducers.

\paragraph{Annotations.}
Reference and coreference annotations over $\Ftwo$ are stored as sparse sorted lists of
critical-cell identifiers, where addition is symmetric difference.  The current core uses a
small inline representation for these lists: the common empty, singleton, and two-label
cases avoid a heap allocation, and longer annotations spill to dynamic storage.  This
optimization is shared by the compact C++ backend, the Python bindings, and the native
GUDHI-view path.  Over an odd prime field $\Fp$, annotations become sparse sorted vectors of pairs
$(\hbox{critical id},\hbox{coefficient})$, and boundary/coboundary signs are reduced
modulo $p$.  The reducer uses an inverse annotation store: for each critical label, it
maintains the list of annotations currently containing that label.  This makes pivot
elimination local to the annotations that actually contain the pivot label.  The
instrumentation records the initial annotation volume, inverse-list entries, candidate
scans, applied updates, total input/output annotation sizes, and maximum annotation sizes.
Composite coefficient rings $\mathbb{Z}_n$ are deliberately rejected by the barcode API;
the current extension is prime-field persistence, not general ring-valued persistence.

\paragraph{Morse sequence strategies.}
We currently compare a strategy portfolio of sequence constructors:
\begin{description}
  \item[\texttt{saturated}] A valid F-sequence constructor that greedily inserts available
        same-level regular pairs, and otherwise inserts a critical simplex.  It pairs until
        it is stuck, marks one fillable simplex critical, then continues.
  \item[\texttt{plateau-greedy}] A variant of \texttt{saturated} that chooses critical
        insertions intended to unlock more same-level pair opportunities.  It pairs until
        it is stuck, marks one strategically chosen fillable simplex critical, then
        continues.  Among the fillable candidates, it scores each candidate by how many
        currently ``almost pairable'' cofaces it would unlock.
  \item[\texttt{same-level-reduction}] An experimental same-level free-face collapse
        strategy.  It discovers free-face pairs inside each filtration level, then emits
        the reverse order as an F-sequence so that the existing reference-map recurrence
        applies.  It collapses until there is no free face, then marks all leftover
        simplices critical.
  \item[\texttt{f-max/min}] Paper-faithful implementations of \textbf{Max}$(S,F)$ and
        \textbf{Min}$(S,F)$ from the Morse-sequence algorithms paper.  The former scans in
        increasing filtration and dimension, while the latter scans in decreasing
        filtration and dimension and prepends the decreasing operations.
  \item[\texttt{flooding-max/min}] Flooding constructors restricted respectively to
        maximal/increasing operations and minimal/decreasing operations.  Candidate
        selection uses extremal priority queues inside each level.
  \item[\texttt{flooding-minmax/maxmin}] Mixed flooding constructors.  The former tries
        minimal/decreasing moves first, while the latter tries maximal/increasing moves
        first.
\end{description}
The interface is intentionally strategy-based, so future Morse sequence algorithms can be
added without changing the persistence entry points.

\paragraph{Reference and coreference paths.}
For an F-sequence, the reference recurrence computes annotations in the direction of
simplex insertion.  For a regular pair $(\sigma,\tau)$ it updates
\[
  R(\tau)=0, \qquad
  R(\sigma)=\bigoplus_{\eta\in \partial \tau,\ \eta\ne\sigma} R(\eta).
\]
We also implemented the dual coreference recurrence.  It scans the sequence backward and
uses coboundaries:
\[
  C(\sigma)=0, \qquad
  C(\tau)=\bigoplus_{\eta\in \mathrm{cobd}(\sigma),\ \eta\ne\tau} C(\eta).
\]
For the \texttt{same-level-reduction} strategy we added a native coreference builder:
it records the free-face collapses, computes coreferences in the natural reverse direction,
and returns the compatible F-sequence for validation and comparison.

\subsection{Benchmark families}

The current benchmark script generates or loads the following families.
\begin{description}
  \item[Lower-star.] Random facet closures with random vertex values.  A simplex receives
        the maximum value of its vertices.
  \item[Plateau.] Random monotone integer-level filtrations over facet closures.  These are
        not forced through a lower-star construction, so same-level plateaus are explicit.
  \item[Rips.] Random distance matrices with Vietoris--Rips filtrations truncated at
        dimension two.
  \item[CAM-style sphere Rips.] Random points on the unit sphere $S^4$ in $\mathbb{R}^5$,
        using the threshold profile from the CAM experiments, downscaled for the prototype.
  \item[Terrain.] Triangulated two-dimensional grids with lower-star filtrations from a
        smooth synthetic scalar field.  The field combines a low-frequency trend,
        sinusoidal oscillations, seeded Gaussian bumps, and small noise, then quantizes the
        vertex values to a small number of levels.  These examples keep spatial structure
        while still producing explicit plateaus.
  \item[Image-grid.] Triangulated two-dimensional grids with stronger image-like plateaus.
        The vertex values come from coarse block regions modified by a seeded rectangle,
        disk, and diagonal stripe, then quantized to a few levels.  These are not cubical
        complexes in the current experiment: the grid is triangulated and then treated as
        an ordinary filtered simplicial complex.  They are intended as a structured
        proxy for grayscale image filtrations.
  \item[Roadmap Rips.] Three distance-matrix datasets from the persistent-homology roadmap
        benchmark collection: random50-16d, random100-4d, and senate104.
\end{description}

For GUDHI comparisons, the GUDHI simplex tree is fed exactly the same simplices and scalar
filtration values stored in our finalized complex.  Thus the comparison is on the same
filtered complex.  The GUDHI timing is used as a current CAM-style baseline, not as a claim
about the original CAM implementation.

\subsection{Correctness checks}

Every benchmark validates the Morse barcode against a standard reduction over the full
filtered complex.  The C++ benchmark path can validate compact barcode signatures without
materializing all intervals in Python.  Additional tests compare:
\begin{itemize}
  \item separate construction of a Morse sequence and reference map;
  \item fused construction of the Morse sequence and reference map;
  \item compact fused reduction input versus a full reference table;
  \item reference reduction versus coreference reduction;
  \item selected examples against GUDHI.
\end{itemize}
Perseus export support is kept as an optional validation path: the benchmark can export the
same finalized complex to \texttt{nmfsimtop}, rank-compress filtration values to positive
integer birth levels, and compare the returned barcode with the standard reducer.  We do
not use Perseus timings in the main experimental comparison because that would mix an
external command-line/file-I/O benchmark with in-process algorithmic timings.
At the time of writing, the C++ test suite, the nanobind build test suite, Python tests, and
GUDHI example validation all pass.

For ratio columns, we report the median with the interquartile interval underneath whenever
the artifact contains several paired observations.  In the native GUDHI tables these
intervals are computed over paired timing repeats.  In the portfolio and synthetic summary
tables whose CSV artifacts already collapse timing repeats, the intervals describe
variation across benchmark cases.  Tables with one aggregate observation per row are
marked as point estimates.

\subsection{Coefficient fields}

The original experimental pipeline was written over $\Ftwo$.  This is still the coefficient
field used for the main tables and for the current GUDHI-facing contribution sketch.
However, the C++ core and Python interface now also expose ordinary persistence,
Morse-reference persistence, and Morse-coreference persistence over prime fields $\Fp$.
The Morse sequence itself is coefficient-independent; the coefficient field enters when
oriented reference/coreference annotations and persistence reductions are computed.

Table~\ref{tab:prime-field-overhead} summarizes a prime-field overhead check on the
synthetic lower-star, plateau, and Rips families.  The run used the
\texttt{saturated}, \texttt{f-max}, \texttt{f-min}, and
\texttt{same-level-reduction} strategies and compared $\mathbb{F}_3$ and
$\mathbb{F}_5$ against the $\Ftwo$ path.  All prime-field
barcodes matched the corresponding standard prime-field reducer.  The ratios below should
not be interpreted as a claim that odd prime fields are faster than $\Ftwo$: they are
best-of-repeat microbenchmarks.  The useful conclusion is that the prime-field path has no
validation mismatches and keeps the same qualitative strategy behavior, with a moderate
overhead in the Morse reducer on these small synthetic cases.

\begin{table}[ht]
\centering
\small
\setlength{\tabcolsep}{3pt}
\begin{tabular}{lrccc}
\toprule
Field & Cases & $\Fp/\Ftwo$ Morse & $\Fp/\Ftwo$ standard
      & Morse speedup vs standard \\
\midrule
$\mathbb{F}_3$ & 180 & \ratioiqr{1.231}{1.190}{1.267} & \ratioiqr{1.040}{1.022}{1.066} & \ratioiqr{1.466}{1.078}{1.898} \\
$\mathbb{F}_5$ & 180 & \ratioiqr{1.216}{1.175}{1.253} & \ratioiqr{1.041}{1.023}{1.062} & \ratioiqr{1.465}{1.122}{1.969} \\
\bottomrule
\end{tabular}
\caption{Prime-field overhead check on synthetic lower-star, plateau, and Rips complexes.
Validation mismatches: 0.  Morse timings exclude sequence construction, which is
coefficient-independent and reported separately in the CSV artifact.  Ratio entries are
case-level medians with $Q_1$--$Q_3$ underneath over the aggregate benchmark rows.}
\label{tab:prime-field-overhead}
\end{table}

\section{Backend and integration comparisons with GUDHI}

We now separate two comparisons that answer different questions.  The backend comparison
uses the compact C++ Morse backend, currently driven through the Python-facing benchmark
pipeline, and compares it with GUDHI on the same filtered complexes.  This is the comparison
where the Morse pipeline is often faster.  The integration comparison starts instead from
an existing \texttt{Gudhi::Simplex\_tree}.  It measures the cost of a possible GUDHI-facing
plugin/interface, including the additional data needed by the Morse sequence, reference map,
and reducer.  After replacing per-simplex heap allocations by inline local storage for
small vertex, boundary, coboundary, and annotation lists, this second path is now
substantially better than compact import.  It is now competitive with GUDHI end-to-end on
structured plateau grids, but not on the flag/Rips-like validation case.  The reducer
kernel is fast; the view and Morse-frame construction stages remain the integration
bottleneck.

\begin{table}[ht]
\centering
\small
\setlength{\tabcolsep}{3pt}
\begin{tabular}{p{0.16\linewidth}p{0.23\linewidth}p{0.23\linewidth}p{0.27\linewidth}}
\toprule
Comparison & Starting object & Timed Morse path & Current conclusion \\
\midrule
Backend & Compact filtered complex & C++ Morse backend & Faster than GUDHI on selected Roadmap averages. \\
Integration & \texttt{Gudhi::Simplex\_tree} & Direct read-only view or import & Direct view is much faster than import; current Release timings are competitive on structured grids but still favor GUDHI on flag/Rips-like input. \\
\bottomrule
\end{tabular}
\caption{The two GUDHI comparisons reported below.  The first measures the Morse backend;
the second measures the present cost of integrating that backend with a GUDHI simplex tree.}
\label{tab:gudhi-comparison-map}
\end{table}

\subsection{Backend comparison on Roadmap datasets}

Table~\ref{tab:roadmap-sat-reduction} compares \texttt{saturated}, \texttt{f-max},
\texttt{f-min}, and \texttt{same-level-reduction} on three objects from the
PH-roadmap benchmark collection:
\texttt{random50-16d}, \texttt{random100-4d}, and \texttt{senate104}.  These are useful
external reference cases because they are not tuned to our implementation and because they
were already used to compare persistent-homology software.  The present experiment should
be read as a targeted comparison on selected Roadmap objects, not as a full reproduction of
the Roadmap benchmark study.

For each object, we build the same truncated Vietoris--Rips filtered complex for our
pipeline and for GUDHI.  The columns report the number of simplices, critical-cell
percentage, total Morse time, GUDHI time, and the GUDHI/Morse point estimate.  A ratio
larger than $1$ means GUDHI is slower than the Morse pipeline.  GUDHI is timed once per complex and
reused across Morse strategies, because it does not depend on the chosen Morse sequence.

\begin{table}[ht]
\centering
\scriptsize
\setlength{\tabcolsep}{2pt}
\begin{tabular}{llrrrrr}
\toprule
Object & Strategy & Simplices & Critical \% & Morse (\us) & GUDHI (\us) & GUDHI/Morse (point) \\
\midrule
\texttt{random50-16d} & \texttt{saturated} & 20,875 & 89.1 & 6,731 & 9,298 & 1.38 \\
\texttt{random50-16d} & \texttt{f-max} & 20,875 & 89.1 & 5,735 & 9,298 & 1.62 \\
\texttt{random50-16d} & \texttt{f-min} & 20,875 & 89.1 & 6,039 & 9,298 & 1.54 \\
\texttt{random50-16d} & \slred & 20,875 & 99.3 & 5,431 & 9,298 & 1.71 \\
\midrule
\texttt{random100-4d} & \texttt{saturated} & 166,750 & 94.2 & 91,902 & 82,858 & 0.90 \\
\texttt{random100-4d} & \texttt{f-max} & 166,750 & 94.2 & 82,689 & 82,858 & 1.00 \\
\texttt{random100-4d} & \texttt{f-min} & 166,750 & 94.2 & 89,447 & 82,858 & 0.93 \\
\texttt{random100-4d} & \slred & 166,750 & 99.8 & 76,253 & 82,858 & 1.09 \\
\midrule
\texttt{senate104} & \texttt{saturated} & 182,207 & 94.4 & 91,024 & 87,125 & 0.96 \\
\texttt{senate104} & \texttt{f-max} & 182,207 & 94.4 & 80,962 & 87,125 & 1.08 \\
\texttt{senate104} & \texttt{f-min} & 182,207 & 94.4 & 92,470 & 87,125 & 0.94 \\
\texttt{senate104} & \slred & 182,207 & 99.9 & 80,891 & 87,125 & 1.08 \\
\bottomrule
\end{tabular}
\caption{Selected PH-roadmap Rips datasets, with GUDHI timed on the same filtered
complexes.  The \texttt{f-max}, \texttt{f-min}, and \texttt{saturated} rows often have
the same critical-cell ratio, but different reducer costs.  The
\texttt{same-level-reduction} strategy leaves many more cells critical, yet remains
competitive on these instances because it performs less annotation mutation work.  Each
row is a single aggregate observation, so the ratio column is a point estimate.}
\label{tab:roadmap-sat-reduction}
\end{table}

The same Roadmap run also covered the full strategy portfolio, including the flooding
variants and plateau-greedy.  Averaged over the three Roadmap objects,
\texttt{same-level-reduction} and \texttt{f-max} were the two fastest strategies:
\texttt{same-level-reduction} averaged $54{,}192~\us$, while \texttt{f-max} averaged
$56{,}462~\us$.  The \texttt{f-min} strategy averaged $62{,}652~\us$, close to
\texttt{saturated} at $63{,}219~\us$.  The older priority-queue flooding variants ranged
from $74{,}674~\us$ for \texttt{plateau-greedy} to $84{,}056~\us$ for
\texttt{flooding-minmax}.

\begin{table}[ht]
\centering
\scriptsize
\setlength{\tabcolsep}{3pt}
\begin{tabular}{lrrcc}
\toprule
Strategy & Critical \% & Morse (\us) & Std/Morse & GUDHI/Morse \\
\midrule
\slred                   & 99.8 & 54,192 & \ratioiqr{4.74}{3.23}{6.40} & \ratioiqr{1.09}{1.08}{1.40} \\
\texttt{f-max}           & 94.0 & 56,462 & \ratioiqr{4.57}{3.16}{6.01} & \ratioiqr{1.08}{1.04}{1.35} \\
\texttt{f-min}           & 94.0 & 62,652 & \ratioiqr{4.17}{2.84}{5.51} & \ratioiqr{0.94}{0.93}{1.24} \\
\texttt{saturated}       & 94.0 & 63,219 & \ratioiqr{3.80}{2.66}{5.22} & \ratioiqr{0.96}{0.93}{1.17} \\
\texttt{plateau-greedy}  & 94.0 & 74,674 & \ratioiqr{2.95}{2.12}{4.33} & \ratioiqr{0.81}{0.79}{0.92} \\
\texttt{flooding-min}    & 94.0 & 75,659 & \ratioiqr{2.58}{1.90}{4.25} & \ratioiqr{0.80}{0.78}{0.85} \\
\texttt{flooding-max}    & 94.0 & 77,831 & \ratioiqr{2.69}{1.93}{4.16} & \ratioiqr{0.78}{0.76}{0.85} \\
\texttt{flooding-maxmin} & 94.0 & 79,268 & \ratioiqr{2.50}{1.85}{3.97} & \ratioiqr{0.75}{0.75}{0.84} \\
\texttt{flooding-minmax} & 94.0 & 84,056 & \ratioiqr{2.66}{1.86}{3.91} & \ratioiqr{0.71}{0.70}{0.82} \\
\bottomrule
\end{tabular}
\caption{Roadmap averages after adding paper-faithful \texttt{f-max} and \texttt{f-min}.
\texttt{f-min} and \texttt{f-max} have the same critical-cell ratio as
\texttt{saturated} on these objects, but their direct Min/Max loops are faster in this
three-repeat freeze.  Ratio entries are case-level medians with $Q_1$--$Q_3$ underneath
over the three Roadmap objects.}
\label{tab:roadmap-fmax-fmin}
\end{table}

\subsection{Integration benchmark from \texorpdfstring{\texttt{Gudhi::Simplex\_tree}}{Gudhi::Simplex\_tree}}

After adding the GUDHI-shaped C++ entry point, we also ran a fully native simplex-tree
benchmark.  This is a different comparison from the Roadmap table above.  The Roadmap
table compares our Python-facing benchmark pipeline with GUDHI on selected external Rips
complexes.  The native benchmark starts from an actual \texttt{Gudhi::Simplex\_tree} and
compares three in-process paths:
\begin{itemize}
  \item GUDHI's persistent cohomology on the input simplex tree;
  \item our Morse pipeline through the direct read-only \texttt{Simplex\_tree} view;
  \item our Morse pipeline after importing the tree into the compact owning complex.
\end{itemize}

% Generated by morseframes/tools/render_native_gudhi_view_table.py.
% Re-run that script instead of editing this table by hand.
\providecommand{\ratioiqr}[3]{\shortstack[c]{#1\\[-0.25ex]{\tiny #2--#3}}}
\begin{table}[ht]
\centering
\scriptsize
\setlength{\tabcolsep}{2.4pt}
\begin{tabular}{llrrrrcc}
\toprule
Case & Strategy & Simp. & Direct (ms) & Import (ms) & GUDHI (ms) & GUDHI/Direct & GUDHI/Reducer \\
\midrule
\texttt{flag-90-r018} & \texttt{f-max} & 952 & 0.28 & 1.48 & 0.34 & \ratioiqr{1.16}{1.12}{1.30} & \ratioiqr{4.15}{3.95}{4.45} \\
\texttt{flag-90-r018} & \texttt{f-min} & 952 & 0.27 & 1.48 & 0.34 & \ratioiqr{1.17}{1.13}{1.34} & \ratioiqr{4.15}{3.99}{4.81} \\
\texttt{flag-90-r018} & \texttt{plateau-greedy} & 952 & 0.40 & 1.61 & 0.34 & \ratioiqr{0.81}{0.78}{0.94} & \ratioiqr{3.93}{3.83}{4.51} \\
\texttt{flag-90-r018} & \slred & 952 & 0.27 & 1.48 & 0.34 & \ratioiqr{1.22}{1.20}{1.31} & \ratioiqr{3.84}{3.73}{4.32} \\
\midrule
\texttt{grid-24x24-plateau} & \texttt{f-max} & 3,267 & 0.61 & 4.86 & 1.27 & \ratioiqr{1.84}{1.77}{1.99} & \ratioiqr{7.14}{6.94}{7.56} \\
\texttt{grid-24x24-plateau} & \texttt{f-min} & 3,267 & 0.62 & 4.73 & 1.27 & \ratioiqr{1.79}{1.73}{1.91} & \ratioiqr{6.85}{6.74}{7.29} \\
\texttt{grid-24x24-plateau} & \texttt{plateau-greedy} & 3,267 & 1.10 & 5.20 & 1.27 & \ratioiqr{1.04}{0.98}{1.09} & \ratioiqr{6.76}{6.44}{7.14} \\
\texttt{grid-24x24-plateau} & \slred & 3,267 & 1.16 & 4.98 & 1.27 & \ratioiqr{1.26}{1.11}{1.30} & \ratioiqr{3.26}{2.66}{3.46} \\
\bottomrule
\end{tabular}
\caption{Native \texttt{Gudhi::Simplex\_tree} quick benchmark.  \texttt{Direct} uses the read-only \texttt{Simplex\_tree} view and our Morse reducer; \texttt{Import} first copies the complex into the compact owning structure.  \texttt{GUDHI} is GUDHI's in-process persistent cohomology on the same \texttt{Simplex\_tree}.  \texttt{GUDHI/Direct}$<1$ means GUDHI is faster end-to-end; \texttt{GUDHI/Reducer}$>1$ means our reducer alone is faster than GUDHI persistence, excluding view construction and Morse input construction.  Ratio columns show the median with $Q_1$--$Q_3$ underneath when repeat-level rows are available.}
\label{tab:native-gudhi-view-quick}
\end{table}


The direct view path is consistently faster than compact import, so the GUDHI interface is
useful.  The quick table is retained as a small smoke test.  After adding inline storage
for the small local lists used by the view and for small reducer annotations, dense
low-dimensional face lookup when vertex handles are compact, rank-based Morse candidate
queues, byte-valued hot state flags, fixed-dimension ordering inside the native view,
level-aware candidate filtering, and cheaper production-mode sequence emission checks, the
larger 30-repeat Release run gives a more useful picture.

% Generated by morseframes/tools/render_native_gudhi_view_table.py.
% Re-run that script instead of editing this table by hand.
\providecommand{\ratioiqr}[3]{\shortstack[c]{#1\\[-0.25ex]{\tiny #2--#3}}}
\begin{table}[ht]
\centering
\scriptsize
\setlength{\tabcolsep}{2.4pt}
\begin{tabular}{llrrrrcc}
\toprule
Case & Strategy & Simp. & Direct (ms) & Import (ms) & GUDHI (ms) & GUDHI/Direct & GUDHI/Reducer \\
\midrule
\texttt{flag-240-r014} & \texttt{f-max} & 7,010 & 1.95 & 11.71 & 2.36 & \ratioiqr{1.20}{1.14}{1.30} & \ratioiqr{4.57}{4.38}{4.90} \\
\texttt{flag-240-r014} & \texttt{f-min} & 7,010 & 1.79 & 11.54 & 2.36 & \ratioiqr{1.30}{1.24}{1.37} & \ratioiqr{4.39}{4.18}{4.73} \\
\texttt{flag-240-r014} & \texttt{plateau-greedy} & 7,010 & 2.82 & 12.67 & 2.36 & \ratioiqr{0.83}{0.79}{0.90} & \ratioiqr{4.34}{3.93}{4.68} \\
\texttt{flag-240-r014} & \slred & 7,010 & 1.54 & 11.34 & 2.36 & \ratioiqr{1.52}{1.47}{1.60} & \ratioiqr{5.44}{5.17}{5.97} \\
\midrule
\texttt{grid-128x128-plateau} & \texttt{f-max} & 97,283 & 24.42 & 144.59 & 37.68 & \ratioiqr{1.57}{1.51}{1.64} & \ratioiqr{6.25}{6.12}{6.42} \\
\texttt{grid-128x128-plateau} & \texttt{f-min} & 97,283 & 24.55 & 144.95 & 37.68 & \ratioiqr{1.57}{1.49}{1.62} & \ratioiqr{5.92}{5.55}{6.17} \\
\texttt{grid-128x128-plateau} & \texttt{plateau-greedy} & 97,283 & 37.48 & 164.03 & 37.68 & \ratioiqr{1.03}{0.99}{1.04} & \ratioiqr{6.12}{5.80}{6.40} \\
\texttt{grid-128x128-plateau} & \slred & 97,283 & 28.42 & 152.43 & 37.68 & \ratioiqr{1.36}{1.31}{1.38} & \ratioiqr{3.64}{3.50}{3.72} \\
\midrule
\texttt{grid-64x64-plateau} & \texttt{f-max} & 24,067 & 3.59 & 31.53 & 6.98 & \ratioiqr{2.03}{1.80}{2.20} & \ratioiqr{154.88}{145.42}{167.98} \\
\texttt{grid-64x64-plateau} & \texttt{f-min} & 24,067 & 3.63 & 31.79 & 6.98 & \ratioiqr{2.00}{1.75}{2.15} & \ratioiqr{144.39}{118.14}{166.78} \\
\texttt{grid-64x64-plateau} & \texttt{plateau-greedy} & 24,067 & 5.75 & 33.78 & 6.98 & \ratioiqr{1.24}{1.17}{1.32} & \ratioiqr{118.38}{107.09}{128.03} \\
\texttt{grid-64x64-plateau} & \slred & 24,067 & 3.68 & 31.45 & 6.98 & \ratioiqr{1.96}{1.75}{2.16} & \ratioiqr{115.29}{102.19}{126.80} \\
\midrule
\texttt{grid-96x96-plateau} & \texttt{f-max} & 54,531 & 12.30 & 79.82 & 20.29 & \ratioiqr{1.78}{1.62}{1.92} & \ratioiqr{7.07}{6.98}{7.45} \\
\texttt{grid-96x96-plateau} & \texttt{f-min} & 54,531 & 12.18 & 80.04 & 20.29 & \ratioiqr{1.77}{1.67}{1.87} & \ratioiqr{7.21}{6.88}{7.56} \\
\texttt{grid-96x96-plateau} & \texttt{plateau-greedy} & 54,531 & 21.03 & 89.46 & 20.29 & \ratioiqr{1.03}{0.98}{1.06} & \ratioiqr{6.90}{6.52}{7.21} \\
\texttt{grid-96x96-plateau} & \slred & 54,531 & 16.88 & 85.07 & 20.29 & \ratioiqr{1.29}{1.20}{1.33} & \ratioiqr{3.15}{3.03}{3.25} \\
\bottomrule
\end{tabular}
\caption{Native \texttt{Gudhi::Simplex\_tree} larger lean benchmark.  \texttt{Direct} uses the read-only \texttt{Simplex\_tree} view and our Morse reducer; \texttt{Import} first copies the complex into the compact owning structure.  \texttt{GUDHI} is GUDHI's in-process persistent cohomology on the same \texttt{Simplex\_tree}.  \texttt{GUDHI/Direct}$<1$ means GUDHI is faster end-to-end; \texttt{GUDHI/Reducer}$>1$ means our reducer alone is faster than GUDHI persistence, excluding view construction and Morse input construction.  Ratio columns show the median with $Q_1$--$Q_3$ underneath when repeat-level rows are available.}
\label{tab:native-gudhi-large-lean-r30}
\end{table}


Table~\ref{tab:native-gudhi-large-lean-r30} is the current integration checkpoint.  On
\texttt{grid-64x64-plateau}, all three reported strategies beat GUDHI end-to-end; on
\texttt{grid-96x96-plateau}, \texttt{f-max} and \texttt{f-min} are essentially at parity;
on \texttt{grid-128x128-plateau}, GUDHI is still faster; and on the flag/Rips-like case
GUDHI remains clearly faster.  The reducer-only ratio remains favorable throughout, so the
result does not argue against the Morse reduction itself.  It identifies view
construction, Morse-sequence construction, and fused reference/reduction-input
construction as the next integration bottlenecks.

% Generated by morseframes/tools/render_native_gudhi_stage_profile.py.
% Re-run that script instead of editing this table by hand.
\begin{table}[ht]
\centering
\scriptsize
\setlength{\tabcolsep}{1.8pt}
\begin{tabular}{llrrrrrrrrl}
\toprule
Case & Strategy & Extract & Bound. & Cobd. & Order & Seq. & Frame & Reducer & Direct & Largest \\
\midrule
\texttt{flag-90-r018} & \texttt{f-max} & 0.04 & 0.03 & 0.04 & 0.01 & 0.07 & 0.07 & 0.10 & 0.37 & \texttt{frame-extra} \\
\texttt{flag-90-r018} & \texttt{f-min} & 0.04 & 0.03 & 0.04 & 0.01 & 0.06 & 0.05 & 0.10 & 0.34 & \texttt{sequence} \\
\texttt{flag-90-r018} & \texttt{plateau-greedy} & 0.04 & 0.03 & 0.04 & 0.01 & 0.24 & 0.04 & 0.11 & 0.51 & \texttt{sequence} \\
\texttt{flag-90-r018} & \texttt{same-level} & 0.04 & 0.03 & 0.04 & 0.01 & 0.05 & 0.06 & 0.11 & 0.34 & \texttt{frame-extra} \\
\midrule
\texttt{grid-24x24-plateau} & \texttt{f-max} & 0.25 & 0.13 & 0.09 & 0.06 & 0.14 & 0.13 & 0.20 & 1.00 & \texttt{extract} \\
\texttt{grid-24x24-plateau} & \texttt{f-min} & 0.25 & 0.13 & 0.09 & 0.06 & 0.16 & 0.21 & 0.20 & 1.10 & \texttt{extract} \\
\texttt{grid-24x24-plateau} & \texttt{plateau-greedy} & 0.25 & 0.13 & 0.09 & 0.06 & 0.70 & 0.10 & 0.20 & 1.54 & \texttt{sequence} \\
\texttt{grid-24x24-plateau} & \texttt{same-level} & 0.25 & 0.13 & 0.09 & 0.06 & 0.21 & 0.28 & 0.55 & 1.58 & \texttt{frame-extra} \\
\bottomrule
\end{tabular}
\caption{Native GUDHI plugin stage profile in milliseconds.  The first four timing columns split construction of the read-only \texttt{Simplex\_tree} view into simplex extraction, boundary lookup, coboundary construction, and filtration ordering.  \texttt{Seq.} times sequence construction alone.  \texttt{Frame extra} is the fused reduction-input time minus sequence-only time, so it estimates reference/reduction-input overhead.}
\label{tab:native-gudhi-stage-profile-quick}
\end{table}


% Generated by morse_persistence/tools/render_native_gudhi_stage_profile.py.
% Re-run that script instead of editing this paragraph by hand.
In this quick native profile, the largest pre-reducer stage is boundary lookup in 1 row, reference/reduction-input overhead in 2 rows, filtration ordering in 1 row, and sequence construction in 4 rows, out of 8 summarized rows.  Averaged over the reported rows, the direct-path time is split approximately as 8.1\% simplex extraction, 12.3\% boundary lookup, 5.4\% coboundary construction, 11.6\% filtration ordering, 20.9\% sequence-only construction, 17.0\% reference/reduction-input overhead, and 25.5\% reducer time.  The GUDHI/direct end-to-end ratio ranges from 0.50 to 0.97, while the main-strategy GUDHI/reducer ratio ranges from 1.59 to 5.72.  This split now shows a more balanced cost distribution: reducer time is no longer the single dominant component, while extraction and sequence construction remain the largest non-reducer costs.


\begin{table}[ht]
\centering
\small
\setlength{\tabcolsep}{3pt}
\begin{tabular}{p{0.30\linewidth}rp{0.50\linewidth}}
\toprule
Direct-view stage group & Share & Interpretation \\
\midrule
Read-only view construction & 37.4\% & Simplex extraction, boundary lookup, coboundaries, and order arrays. \\
Morse frame construction & 37.9\% & Morse sequence plus reference/reduction-input construction. \\
Reducer kernel & 25.5\% & Annotation reduction after the Morse input has been built. \\
\midrule
All pre-reducer work & 74.5\% & Current native-plugin bottleneck before the reducer starts. \\
\bottomrule
\end{tabular}
\caption{Compact decomposition of the direct native \texttt{Simplex\_tree} view path,
averaged over the quick plugin stage-profile rows.  The percentages are rounded from the
generated stage profile and therefore need not sum to exactly $100\%$.}
\label{tab:native-gudhi-overhead-summary}
\end{table}


Table~\ref{tab:native-gudhi-overhead-summary} is the key integration diagnostic.  Inline
local storage reduced view construction substantially, especially coboundary construction,
and the dense low-dimensional face lookup reduces boundary construction on compact-handle
cases.  Rank-based Morse candidate queues and one-pass rank/level/dimension setup reduce
sequence-construction overhead in the important \texttt{f-max}/\texttt{f-min} rows.
Byte-valued hot state flags remove \texttt{std::vector<bool>} proxy overhead in the
sequence/reference construction loops.
Inline storage for small annotations and the specialized \(\mathbb{Z}_2\) reducer update
also reduced the reducer share.  Roughly three quarters of the direct native path is
still spent before the reducer starts.
This explains both the improvement and the remaining below-parity native rows: the compact
backend avoids most of this \texttt{Simplex\_tree}-view construction cost, while the
GUDHI-facing path must rebuild enough incidence and ordering information for the Morse
frame.  The current evidence is therefore ``competitive on structured grids, still
unresolved on flag/Rips-like inputs'', not a uniform speedup claim.

The important observation is counterintuitive: \texttt{same-level-reduction} leaves almost every
cell critical on these datasets, but it still improves running time.  For example, on
Roadmap 100, \texttt{saturated} has $157{,}154$ critical cells and \texttt{same-level-reduction} has
$166{,}470$, yet \texttt{same-level-reduction} is faster.  The instrumentation explains this:

\begin{table}[ht]
\centering
\small
\begin{tabular}{lrrrrr}
\toprule
Strategy & Initial labels & Pivots & XOR input & XOR output & Morse (\us) \\
\midrule
saturated   & 208,954 & 152   & 89,676 & 50,040 & 89,505 \\
\slred & 166,822 & 4,810 & 10,786 & 420    & 56,289 \\
\bottomrule
\end{tabular}
\caption{Roadmap 100 reducer metrics.  The \slred strategy uses more pivots, but much
smaller annotation updates; the critical-cell ratio alone is therefore misleading.}
\label{tab:roadmap100-metrics}
\end{table}

\section{Strategy selection and profiling}

The Roadmap experiment suggests that no single Morse sequence strategy should be hard-coded
as the default.  We therefore added profiling tools that construct candidate Morse frames,
record cheap work proxies, and compare the profile-selected strategy with the actually
fastest measured strategy.  This section separates three related questions:
\begin{itemize}
  \item which cheap score should rank strategies;
  \item which subset of strategies should be profiled;
  \item whether structured families, especially image-like grids, need a special policy.
\end{itemize}

\subsection{Roadmap profile selection}
We then tested the new cheap profile mode on the same Roadmap objects.  For each strategy,
the profiler constructs the Morse sequence and compact reference-reduction input, but does
not run pivot elimination.  It records the working set, initial annotation volume, boundary
annotation work, and an estimated reducer-work score.  Table~\ref{tab:roadmap-profile}
compares the profile-selected strategy with the strategy selected by full timings.

\begin{table}[ht]
\centering
\scriptsize
\setlength{\tabcolsep}{3pt}
\begin{tabular}{lllrrrr}
\toprule
Object & Profile best & Timing best & Penalty \% & GUDHI (\us) & Profile portfolio (ms) & Measured portfolio (ms) \\
\midrule
\texttt{random50-16d}  & \slred & \texttt{f-max} & 0.13 & 9,367  & 109.3 & 80.5 \\
\texttt{random100-4d} & \slred & \slred         & 0.00 & 85,158 & 1,160.9 & 887.3 \\
\texttt{senate104}    & \slred & \slred         & 0.00 & 90,015 & 1,109.9 & 922.8 \\
\bottomrule
\end{tabular}
\caption{Profile-based strategy selection on selected PH-roadmap objects using the
time-adjusted \texttt{profile-total-work} score.  On \texttt{random50-16d}, the measured
best strategy is \texttt{f-max}, but the \texttt{same-level-reduction} penalty is only
$0.13\%$.  The portfolio timings are sums over the nine candidate strategies, not timings
for the selected strategy alone.}
\label{tab:roadmap-profile}
\end{table}

This is encouraging, but it also shows a practical limitation of the current prototype:
profiling avoids pivot elimination and barcode construction, yet profiling the full
strategy portfolio is not always cheap enough on large Rips complexes.  The selector should
therefore either profile a smaller candidate set or make the profile path cheaper.

\subsection{Profile-versus-measured calibration}
We added a benchmark mode that compares the profile-selected strategy with the actually
fastest measured Morse strategy on each generated complex.  The output reports the profile
choice, the measured best choice, the measured rank of the profile choice, the timing
penalty, and diagnostic gaps in constructor time and measured reducer work.  We used it to
compare the old \texttt{estimated-reducer-work} score with a new
\texttt{profile-total-work} score, defined as estimated reducer work plus profiled frame
time converted to nanoseconds.

\begin{table}[ht]
\centering
\scriptsize
\setlength{\tabcolsep}{3pt}
\begin{tabular}{llrrrrr}
\toprule
Dataset & Score & Cases & Match \% & Avg. pen. \% & Max pen. \% & Avg. rank \\
\midrule
base families & estimated work & 27 & 22.2 & 39.9 & 190.3 & 3.67 \\
base families & total work     & 27 & 51.9 & 7.4  & 38.9  & 1.89 \\
\midrule
CAM $S^4$ Rips & estimated work & 9 & 44.4 & 3.2 & 15.9 & 1.89 \\
CAM $S^4$ Rips & total work     & 9 & 11.1 & 6.5 & 38.0 & 2.78 \\
\midrule
Roadmap Rips & estimated work & 3 & 66.7 & 2.3 & 6.8 & 1.67 \\
Roadmap Rips & total work     & 3 & 66.7 & 0.0 & 0.1 & 1.33 \\
\bottomrule
\end{tabular}
\caption{Profile-versus-measured calibration with the full current strategy portfolio.
The base-family experiment uses lower-star, plateau, and Rips examples at sizes $16$,
$24$, and $32$ with three seeds.  The CAM experiment uses the downscaled sphere-Rips
family, also at sizes $16$, $24$, and $32$.  The Roadmap experiment uses the three selected
PH-roadmap Rips objects.}
\label{tab:profile-calibration-base}
\end{table}

This mixed result is useful.  The time-adjusted score largely fixes the base-family
failure mode, where the old score often chose a strategy with an attractive reduction input
but too much construction overhead.  However, it is worse on the CAM-style sphere examples,
where measured times are very close and the additional profile-time term perturbs the
ranking.  The selector should therefore become family- or scale-aware rather than use a
single global score.

\subsection{Candidate-gate pilot}
We then added experimental gates that profile only a small subset of the full strategy
portfolio, while still measuring the full portfolio for evaluation.  The
\texttt{family-aware} gate uses the benchmark family label.  The \texttt{shape-aware} gate
uses only cheap complex-shape features recorded in the CSV: dimension counts, edge and
triangle density, boundary/coboundary fan-in, and filtration-level concentration.
Table~\ref{tab:profile-gate} compares these gates with full-portfolio profiling, using the
\texttt{profile-total-work} score in all cases.

\begin{table}[ht]
\centering
\scriptsize
\setlength{\tabcolsep}{3pt}
\begin{tabular}{llrrrrr}
\toprule
Dataset & Gate & Profiled & Match \% & Avg. pen. \% & Max pen. \% & Profile overhead \% \\
\midrule
base families & full portfolio & 9.0 & 51.9 & 7.4 & 38.9 & 68.0 \\
base families & family-aware   & 3.0 & 40.7 & 6.5 & 26.6 & 15.9 \\
base families & shape-aware    & 3.8 & 29.6 & 12.2 & 109.8 & 18.8 \\
\midrule
CAM $S^4$ Rips & full portfolio & 9.0 & 11.1 & 6.5 & 38.0 & 51.9 \\
CAM $S^4$ Rips & family-aware   & 3.0 & 33.3 & 5.9 & 26.2 & 11.1 \\
CAM $S^4$ Rips & shape-aware    & 4.0 & 66.7 & 2.0 & 9.2 & 14.7 \\
\midrule
Roadmap Rips & full portfolio & 9.0 & 66.7 & 0.0 & 0.1 & 125.9 \\
Roadmap Rips & family-aware   & 2.0 & 66.7 & 0.7 & 2.0 & 23.2 \\
Roadmap Rips & shape-aware    & 2.0 & 33.3 & 5.4 & 15.4 & 18.1 \\
\bottomrule
\end{tabular}
\caption{Candidate-gate pilot.  ``Profile overhead'' is the total profile time divided by
the total measured Morse-portfolio time.  The family-aware gate sharply reduces profiling
cost with small average penalties.  The first shape-aware gate helps on the CAM-style
sphere-Rips examples, but is not robust on the base and Roadmap samples.}
\label{tab:profile-gate}
\end{table}

These gates are not final selectors.  In this first pilot, the hand-written
\texttt{family-aware} gate is the most useful simple control: it reduces profiling cost
substantially and keeps penalties small.  The \texttt{shape-aware} gate is useful
diagnostically because it identifies sparse CAM-style sphere-Rips examples, but its
hand-written rules overreact on some small base examples and on Roadmap objects where
\texttt{f-max} and \texttt{same-level-reduction} are close.  The shape fields are
therefore worth keeping, but the rule should be learned or tabulated from more calibration
data rather than tuned by hand.  The hold-out paragraph below is the first step in that
direction.

\subsubsection*{Tabulated candidate gates.}
We added \texttt{calibrate\_profile\_gate.py} to turn calibration CSVs into candidate
tables.  The utility groups \texttt{profile-vs-measured} rows by family and/or cheap shape
bins, then reports which measured winners should be retained in each gate.  On the current
calibration rows, a two-candidate family-level table covers all observed winners:
\begin{itemize}
  \item lower-star: \texttt{f-max}/\texttt{saturated};
  \item plateau and synthetic Rips: \texttt{f-max}/\texttt{f-min};
  \item CAM-style sphere Rips: \texttt{f-min}/\texttt{same-level-reduction};
  \item Roadmap Rips: \texttt{f-max}/\texttt{same-level-reduction}.
\end{itemize}
This is a calibration hint, not a final selector.  The benchmark driver can consume such a
table with \texttt{--profile-candidate-gate tabulated} and
\texttt{--profile-candidate-table}, separating calibration runs from hold-out validation
runs.  The calibration utility also supports \texttt{--include-candidates}, so conservative
tables can retain guard strategies even when they were not winners in the calibration
subset.

\subsubsection*{Hold-out behavior.}
A small fresh base-family hold-out on sizes $16$, $24$, and $32$ with seeds $8$ and $9$
compared the current \texttt{family-aware} gate with an updated conservative tabulated
gate.  Both gates had zero unavailable measured winners.  The conservative tabulated gate
had lower aggregate average penalty, $4.5\%$ versus $6.9\%$, and lower aggregate profiling
overhead, $12.4\%$ versus $15.0\%$.

The tradeoff is concentrated in the plateau family: retaining \texttt{f-max},
\texttt{f-min}, \texttt{saturated}, and \texttt{plateau-greedy} prevents excluded winners
but raises plateau profiling overhead and did not improve the plateau selector on this
small split.  A family-size tabulated gate reduced aggregate profiling overhead to
$10.7\%$, but introduced one unavailable lower-star winner and increased plateau average
penalty to $9.6\%$.  Thus finer table keys are useful diagnostically, but they do not by
themselves solve the plateau selector problem.

\subsection{Fair metric validation and selector audit}
We then compared fixed \texttt{profile-total-work}, fixed \texttt{critical-ratio}, and two
adaptive profile metrics.  The conservative \texttt{adaptive} metric uses
\texttt{critical-ratio} except on Roadmap Rips, where it keeps
\texttt{profile-total-work}.  The structured-grid variant,
\texttt{adaptive-structured}, keeps the Roadmap rule, uses
\texttt{profile-total-work} for image-grid filtrations, and leaves generic plateau and
terrain cases on \texttt{critical-ratio}.  The first version of this comparison used
separate benchmark runs for each metric, which made the conclusion vulnerable to timing
noise.  We therefore added a fair multi-metric mode: a single
\texttt{profile-vs-measured} run can now evaluate several profile metrics against the same
profiles and the same measured portfolio timings.  The generated rows record both the
requested profile metric and the effective metric used after resolving adaptive policies.
They also record cheap shape and plateau descriptors, including level concentration,
entropy, largest-level size, and simplex-to-vertex ratio, so future selector experiments
can be analyzed without rerunning the measured portfolio.

% Generated by morseframes/tools/run_fair_profile_validation.py.
% Re-run that script instead of editing this table by hand.
\begin{table}[ht]
\centering
\small
\begin{tabular}{llrrrrrr}
\toprule
Split & Metric (eff.) & Cases & Match \% & Unavail. & Avg. pen. \% & Max pen. \% & Rank \\
\midrule
Base & PTW (PTW) & 27 & 59.3 & 0 & 4.37 & 39.66 & 1.52 \\
Base & CR (CR) & 27 & 44.4 & 0 & 4.42 & 24.92 & 1.70 \\
Base & Adapt (CR) & 27 & 44.4 & 0 & 4.42 & 24.92 & 1.70 \\
Base & Adapt-S (CR) & 27 & 44.4 & 0 & 4.42 & 24.92 & 1.70 \\
\midrule
CAM $S^4$ & PTW (PTW) & 9 & 66.7 & 0 & 1.29 & 5.43 & 1.44 \\
CAM $S^4$ & CR (CR) & 9 & 66.7 & 0 & 1.29 & 5.43 & 1.44 \\
CAM $S^4$ & Adapt (CR) & 9 & 66.7 & 0 & 1.29 & 5.43 & 1.44 \\
CAM $S^4$ & Adapt-S (CR) & 9 & 66.7 & 0 & 1.29 & 5.43 & 1.44 \\
\midrule
Roadmap & PTW (PTW) & 3 & 100.0 & 0 & 0.00 & 0.00 & 1.00 \\
Roadmap & CR (CR) & 3 & 0.0 & 0 & 57.83 & 69.83 & 3.00 \\
Roadmap & Adapt (PTW) & 3 & 100.0 & 0 & 0.00 & 0.00 & 1.00 \\
Roadmap & Adapt-S (PTW) & 3 & 100.0 & 0 & 0.00 & 0.00 & 1.00 \\
\bottomrule
\end{tabular}
\caption{Fair multi-metric validation of profile-selection scores with the current
\texttt{family-aware} candidate gate.  PTW denotes \texttt{profile-total-work}, CR denotes
\texttt{critical-ratio}, Adapt denotes \texttt{adaptive}, and Adapt-S denotes
\texttt{adaptive-structured}; the metric in parentheses is
the effective metric after resolving the adaptive policy.  Within each split, all
requested metrics are evaluated against the same profiles and the same measured portfolio
timings.  ``Unavail.'' counts cases where the measured fastest strategy was not in the
profiled candidate set.  ``Rank'' is the average measured-time rank of the
profile-selected strategy.}
\label{tab:fair-profile-metrics}
\end{table}


% Generated by morse_persistence/tools/run_fair_profile_validation.py.
% Re-run that script instead of editing this paragraph by hand.
On this validation split with five timing repeats, the controlled comparison reduces the base-family aggregate average penalty from $8.34\%$ for \texttt{profile-total-work} to $2.22\%$ for \texttt{critical-ratio}/\texttt{adaptive}.
For CAM-style Rips all three metrics select the same strategies on this split, with $0.66\%$ average penalty.
For Roadmap Rips, the average penalties are \texttt{profile-total-work} $2.31\%$, \texttt{critical-ratio} $3.72\%$, \texttt{adaptive} $2.31\%$.
Thus adaptive scoring with the current \texttt{family-aware} gate remains the most promising experimental selector, and future selector comparisons should use the multi-metric mode rather than independent timing runs.


After generating the report split, the extended hold-out, the plateau-focused hold-out,
the terrain hold-out, the image-grid hold-out, and the larger image-grid stress split, we
combine their aggregate rows into the selector-decision audit in
Table~\ref{tab:selector-decision-summary}.

% Generated by morseframes/tools/summarize_selector_decisions.py.
% Re-run that script instead of editing this table by hand.
\begin{table}[ht]
\centering
\scriptsize
\setlength{\tabcolsep}{3pt}
\begin{tabular}{@{}llrlrrl@{}}
\toprule
Source & Split & Cases & Best metric(s) & Cur./best pen. \% & Gap pp & Decision \\
\midrule
Report & Base & 27 & Adapt/Adapt-S/CR & 2.22/2.22 & 0.00 & keep \\
Report & CAM & 9 & Adapt/Adapt-S/CR/PTW & 0.66/0.66 & 0.00 & keep \\
Report & Roadmap & 3 & Adapt/Adapt-S/PTW & 2.31/2.31 & 0.00 & keep \\
Extended & Base hold-out & 60 & Adapt/Adapt-S/CR & 3.68/3.68 & 0.00 & keep \\
Extended & CAM hold-out & 9 & Adapt/Adapt-S/CR/PTW & 2.17/2.17 & 0.00 & keep \\
Extended & Roadmap hold-out & 3 & Adapt/Adapt-S/PTW & 0.99/0.99 & 0.00 & keep \\
Plateau & Plateau hold-out & 100 & Adapt/Adapt-S/CR/ERW & 5.05/5.05 & 0.00 & keep \\
Terrain & Terrain hold-out & 20 & Adapt/Adapt-S/CR/ERW & 1.74/1.74 & 0.00 & keep \\
Image & Image hold-out & 20 & Adapt-S/Psec/PTW & 2.86/1.70 & 1.16 & switch candidate \\
Image+ & Image stress & 20 & Adapt-S/Psec/PTW & 0.49/0.42 & 0.07 & watch best \\
\bottomrule
\end{tabular}
\caption{Selector-decision audit over the aggregate rows of the fair
profile-selection validation summaries.  The current metric is
\texttt{adaptive}.  A zero gap means that the current selector is tied with
the best observed average penalty for that validation scope.}
\label{tab:selector-decision-summary}
\end{table}


% Generated by morse_persistence/tools/summarize_selector_decisions.py.
% Re-run that script instead of editing this paragraph by hand.
The selector-decision audit keeps the current \texttt{adaptive} metric in 8 of 10 aggregate validation scopes.
The following scopes have close alternatives to watch: image-stress (Adapt-S/Psec/PTW, 0.07 pp).
The following scopes exceed the promotion margin: image-holdout (Adapt-S/Psec/PTW, 1.16 pp).


Table~\ref{tab:selector-decision-summary} treats \texttt{adaptive} as the current control.
Under that convention, the only nonzero gap is the image-grid branch: the image hold-out is
a switch candidate and the larger image stress split is a watch case.  If the same audit is
rerun with \texttt{adaptive-structured} as the current metric, all ten aggregate validation
scopes are keep/tie decisions.  This does not yet prove that the structured branch should
become the public default, because the image-grid evidence is still only two aggregate
splits, but it does change the interpretation: \texttt{adaptive-structured} is now the main
selector candidate to validate next, not merely an exploratory diagnostic.

\subsection{Image-grid selector diagnostic}

To make the image-grid question more concrete, we also compare requested
\texttt{critical-ratio} against the best requested total-work/time metric on each
validation case, then summarize the margin against the cheap shape descriptors now stored
in the profile CSVs.

% Generated by morseframes/tools/analyze_selector_features.py.
% Re-run that script instead of editing this table by hand.
\begin{table}[ht]
\centering
\small
\begin{tabular}{llrrrrrr}
\toprule
Split & Family & Cases & TW wins & CR wins & Ties & CR pen. & TW pen. \\
\midrule
Image hold-out & image-grid & 20 & 0 & 3 & 17 & 1.36 & 2.02 \\
Image stress & image-grid & 20 & 3 & 2 & 15 & 4.57 & 4.11 \\
Plateau hold-out & plateau & 100 & 24 & 4 & 72 & 4.21 & 1.25 \\
Terrain hold-out & terrain & 20 & 1 & 1 & 18 & 6.64 & 6.53 \\
All & all & 271 & 65 & 22 & 184 & 5.92 & 2.34 \\
\bottomrule
\end{tabular}
\caption{Selector-feature diagnostic comparing requested \texttt{critical-ratio} against the best requested total-work/time metric on each validation case.  Positive evidence for the structured branch appears when the TW penalty is lower than the CR penalty.}
\label{tab:selector-feature-diagnostic}
\end{table}


% Generated by morse_persistence/tools/analyze_selector_features.py.
% Re-run that script instead of editing this paragraph by hand.
The selector-feature diagnostic compares requested \texttt{critical-ratio} with the best requested total-work/time metric on 271 validation cases.  The aggregate average penalties are $3.45\%$ for critical-ratio and $4.51\%$ for total-work/time.
For image-grid filtrations, the total-work/time branch is best on the small hold-out by 1.16 percentage points and on the larger stress split by 0.07 percentage points.  This supports keeping the structured branch as a diagnostic candidate while requiring more validation before promoting it.
The regenerated image rows now expose descriptor coverage for this comparison: the average simplex-to-vertex ratio rises from 5.10 on the hold-out to 5.59 on the stress split, which gives the next calibration pass a concrete scale variable to test.
The terrain hold-out remains a useful control because it has grid structure but does not show the same large total-work/time margin (-0.36 percentage points).


Thus the image-like branch is a useful diagnostic and a credible promotion candidate, but
not yet a default API policy.  The next calibration pass should regenerate all validation
splits with the new descriptor columns and then test whether a stable threshold in scale,
simplex-to-vertex ratio, or plateau concentration predicts the cases where
total-work/time scoring is worth using.

\section{Synthetic strategy comparisons}

\subsection{Synthetic scale experiments}

Table~\ref{tab:synthetic-scale} reports representative scale results for size $48$ in
the synthetic families.  Timing values are averaged over three seeds in the benchmark
summary, while the ratio column reports the median with $Q_1$--$Q_3$ underneath across
those seeds.
The table includes \texttt{f-max} and \texttt{f-min}; earlier drafts of this note only
showed \texttt{saturated} and \slred in this section, which made the synthetic scale
comparison lag behind the current strategy portfolio.

% Generated by morseframes/tools/render_synthetic_scale_table.py.
% Re-run that script instead of editing this table by hand.
\providecommand{\ratioiqr}[3]{\shortstack[c]{#1\\[-0.25ex]{\tiny #2--#3}}}
\begin{table}[ht]
\centering
\small
\setlength{\tabcolsep}{3pt}
\begin{tabular}{llrrrrc}
\toprule
Family & Strategy & Cases & Avg. simplices & Critical \% & Morse (\us) & Std/Morse \\
\midrule
lower-star & \texttt{f-max} & 3 & 317.7 & 18.8 & 48.4 & \ratioiqr{1.83}{1.82}{1.92} \\
lower-star & \texttt{f-min} & 3 & 317.7 & 18.8 & 56.0 & \ratioiqr{1.61}{1.60}{1.64} \\
lower-star & \texttt{saturated} & 3 & 317.7 & 18.8 & 90.2 & \ratioiqr{1.01}{0.97}{1.10} \\
lower-star & \slred & 3 & 317.7 & 36.2 & 68.9 & \ratioiqr{1.24}{1.24}{1.34} \\
\midrule
plateau & \texttt{f-max} & 3 & 317.7 & 22.1 & 52.8 & \ratioiqr{1.68}{1.57}{1.78} \\
plateau & \texttt{f-min} & 3 & 317.7 & 22.1 & 59.4 & \ratioiqr{1.51}{1.46}{1.52} \\
plateau & \texttt{saturated} & 3 & 317.7 & 22.1 & 74.6 & \ratioiqr{1.19}{1.17}{1.23} \\
plateau & \slred & 3 & 317.7 & 38.5 & 62.4 & \ratioiqr{1.39}{1.34}{1.44} \\
\midrule
rips & \texttt{f-max} & 3 & 2,770.3 & 67.3 & 755.8 & \ratioiqr{4.17}{3.92}{4.65} \\
rips & \texttt{f-min} & 3 & 2,770.3 & 67.3 & 747.7 & \ratioiqr{4.58}{4.21}{4.66} \\
rips & \texttt{saturated} & 3 & 2,770.3 & 67.3 & 1,065.3 & \ratioiqr{3.17}{2.97}{3.38} \\
rips & \slred & 3 & 2,770.3 & 94.2 & 664.5 & \ratioiqr{5.46}{4.63}{5.63} \\
\bottomrule
\end{tabular}
\caption{Representative synthetic scale results at size $48$ for the core same-level strategy family.  On small lower-star and plateau examples, the Morse machinery is slower than standard persistence because overhead dominates.  On larger Rips examples, Morse reduction becomes faster.  The \texttt{Std/Morse} ratio shows the median with $Q_1$--$Q_3$ underneath over the listed cases.}
\label{tab:synthetic-scale}
\end{table}


% Generated by morse_persistence/tools/render_synthetic_scale_table.py.
% Re-run that script instead of editing this paragraph by hand.
These experiments suggest that the present implementation is not yet optimized for small inputs.  They also show why \texttt{f-max} and \texttt{f-min} must remain in the measured portfolio: the fastest core strategy in this representative run is lower-star: \texttt{f-max}, plateau: \texttt{f-max}, rips: \texttt{f-max}.
On denser Rips complexes, several Morse strategies begin to beat standard persistence, which is the regime where standard reduction becomes expensive.


\subsection{Flooding sequence pilot}

We then added the flooding family to the measured portfolio, including both mixed
directions.  The current version uses priority queues inside each filtration level:
maximal queues prefer higher-dimensional cells later in the level order, while minimal
queues prefer lower-dimensional cells earlier in the level order.  Table~\ref{tab:flooding-pilot}
reports a focused pilot on the lower-star, plateau, and Rips families with sizes $16$ and
$24$, two seeds, and three timing repeats.  GUDHI was timed on the same filtered
complexes.

\begin{table}[ht]
\centering
\small
\setlength{\tabcolsep}{3pt}
\begin{tabular}{lrrrrc}
\toprule
Strategy & Cases & Avg. simplices & Critical \% & Morse (\us) & GUDHI/Morse \\
\midrule
\texttt{f-max}           & 12 & 171.1 & 34.0 & 46.1 & \ratioiqr{1.17}{0.76}{1.41} \\
\texttt{f-min}           & 12 & 171.1 & 34.0 & 47.9 & \ratioiqr{1.03}{0.76}{1.30} \\
\texttt{saturated}       & 12 & 171.1 & 34.0 & 56.0 & \ratioiqr{0.87}{0.61}{1.04} \\
\slred                   & 12 & 171.1 & 60.3 & 59.2 & \ratioiqr{0.76}{0.64}{0.87} \\
\texttt{flooding-min}    & 12 & 171.1 & 34.0 & 74.4 & \ratioiqr{0.59}{0.42}{0.63} \\
\texttt{plateau-greedy}  & 12 & 171.1 & 34.0 & 75.3 & \ratioiqr{0.55}{0.52}{0.67} \\
\texttt{flooding-max}    & 12 & 171.1 & 34.0 & 76.6 & \ratioiqr{0.53}{0.49}{0.61} \\
\texttt{flooding-minmax} & 12 & 171.1 & 34.0 & 77.5 & \ratioiqr{0.55}{0.45}{0.60} \\
\texttt{flooding-maxmin} & 12 & 171.1 & 34.0 & 78.0 & \ratioiqr{0.53}{0.49}{0.57} \\
\bottomrule
\end{tabular}
\caption{Base-family benchmark after adding paper-faithful \texttt{f-max} and
\texttt{f-min}.  These two strategies produce the same critical-cell ratio as
\texttt{saturated} and the flooding variants on these examples, but avoid the heavier
priority-queue overhead of the earlier flooding implementation.  The GUDHI/Morse column is
the case-level median with $Q_1$--$Q_3$ underneath over the 12 benchmark cases.}
\label{tab:flooding-pilot}
\end{table}

The paired comparison between \texttt{flooding-minmax} and
\texttt{flooding-maxmin} remains diagnostically useful: the distinction is real at the
sequence level, but the first extremal priority policy did not create a substantially
different Morse complex on these benchmark families.  The new \texttt{f-max} and
\texttt{f-min} implementations clarify the situation.  The paper-faithful Max/Min loops
give the same critical-cell ratio as the flooding variants here, but with much lower
constructor overhead.

\section{Coreference experiment}

Because CAM is cohomological in spirit, we also implemented a coreference path.  The goal
was to test whether the natural dual of same-level reduction would behave more like a cohomological
algorithm.  The result was negative for the current data structure and recurrence.

Table~\ref{tab:coreference-roadmap} compares the two directions for \texttt{same-level-reduction}
on the Roadmap datasets.  The reference direction uses the usual boundary-based recurrence;
the coreference direction uses the coboundary-based recurrence.

\begin{table}[ht]
\centering
\scriptsize
\setlength{\tabcolsep}{2pt}
\begin{tabular}{llrrrrr}
\toprule
Dataset & Direction & Total (\us) & Std/Morse (point) & Workset & Initial labels & XOR input \\
\midrule
Roadmap 50  & reference   & 5,295    & 5.32 & 20,799  & 20,916  & 2,788 \\
Roadmap 50  & coreference & 44,186   & 0.64 & 20,870  & 25,880  & 764,215 \\
\midrule
Roadmap 100 & reference   & 54,043   & 11.46 & 166,610 & 166,822 & 10,786 \\
Roadmap 100 & coreference & 1,430,049 & 0.43 & 166,747 & 186,169 & 17,732,185 \\
\midrule
Roadmap 104 & reference   & 60,462   & 2.39 & 182,092 & 182,267 & 11,117 \\
Roadmap 104 & coreference & 564,029  & 0.26 & 182,203 & 199,965 & 3,307,174 \\
\bottomrule
\end{tabular}
\caption{Reference versus coreference reduction for the \texttt{same-level-reduction} sequence.
The coreference path is correct, but its annotation update volume explodes.  Each ratio is
a single aggregate observation for the corresponding dataset and direction.}
\label{tab:coreference-roadmap}
\end{table}

The reason is structural.  Boundaries are small and dimension-bounded: in a two-dimensional
complex an edge has two boundary vertices and a triangle has three boundary edges.
Coboundaries are not bounded in this way.  In a dense Rips complex, a vertex or an edge can
have many cofaces.  Moreover, a simplex that is free among the currently active same-level
cofaces need not be globally free in the whole filtered complex.  Higher-level cofaces may
already carry nontrivial annotations in the coreference recurrence.  Thus the dual update
\[
  C(\tau)=\bigoplus_{\eta\in \mathrm{cobd}(\sigma),\ \eta\ne\tau} C(\eta)
\]
can accumulate very large annotations.  On Roadmap 100, the total XOR input increases from
$10{,}786$ in the reference direction to $17{,}732{,}185$ in the coreference direction.

This should not be interpreted as evidence against CAM.  CAM is a compressed annotation
matrix algorithm for persistent cohomology: it processes inserted simplices and computes
annotations of their boundaries.  It is dual/cohomological, but it is not the same as the
naive coboundary-based coreference recurrence above.  Our coreference experiment shows that
this particular dual Morse recurrence is a poor match for dense Rips coboundaries.

\section{Interpretation and next steps}

The experiments point to five conclusions.

\subsection{Critical-cell ratio is not sufficient}
A sequence with fewer critical cells is not necessarily faster.  The Roadmap experiments
show that \texttt{same-level-reduction} can be faster despite making almost every simplex
critical.  The fair multi-metric validation also shows the other side of the same issue:
\texttt{critical-ratio} is a good cheap selector on the base synthetic families in the
current split, but it fails on the Roadmap Rips objects.  The useful predictors are
therefore family-dependent.  On dense Roadmap-style Rips complexes, total profiled work
and annotation volume are more reliable than the critical-cell count alone.

\subsection{Strategy profiling is a useful selection layer}
The Python interface now includes both a measured portfolio mode and a cheaper profile
mode.  The profile mode constructs candidate Morse frames and records cheap shape and
reducer-work features before full timing.  The current practical control is the
\texttt{family-aware} candidate gate combined with the experimental \texttt{adaptive}
profile metric: it uses \texttt{critical-ratio} except on Roadmap Rips, where it uses
\texttt{profile-total-work}.  Table~\ref{tab:fair-profile-metrics} is now
generated directly from the validation CSVs, and the accompanying generated paragraph is
derived from the same rows.  The selector-decision audit then combines the report,
extended hold-out, plateau-focused hold-out, terrain hold-out, image-grid hold-out, and
image-grid stress summaries into a single keep/watch/switch table.  This makes the result reproducible and
prevents the report from drifting away from the benchmark artifacts.  It is still a
selector experiment, not a final theorem: the image-grid splits currently favor the
structured adaptive branch, and the alternative audit says that \texttt{adaptive-structured}
would remove the only current aggregate selector gap.  The right conclusion is therefore
stronger than in the previous draft but still cautious: keep \texttt{adaptive} as the
control, and validate \texttt{adaptive-structured} as the next candidate before making it a
user-facing default.

\subsection{Native GUDHI integration is useful and now partly competitive}
The direct \texttt{Gudhi::Simplex\_tree} view has changed from a diagnostic prototype into
a credible integration path.  It avoids the large cost of importing the tree into the
compact owning complex, and the reducer kernel is faster than GUDHI persistence once the
Morse input has been built.  In the current Release native benchmark, the direct path is
faster than GUDHI on the smaller structured grid and essentially tied on the medium grid,
but GUDHI remains faster on the flag/Rips-like case and on the largest grid.  The
integration claim should therefore be phrased as ``competitive on structured grids, not
yet uniformly faster'', rather than as a global speedup claim.  The remaining engineering
bottleneck is not only the reducer: a large share of the direct native path is still spent
before reduction, in view construction and Morse-frame construction.

\subsection{Reference reduction is the useful direction for now}
The native coreference path is correct, but it is much slower on the Rips-style datasets
because coboundary fan-in causes large annotation diffusion.  The promising path is
currently to construct good Morse sequences and then use the reference reduction, not the
naive coboundary-based coreference recurrence.  This does not contradict CAM: CAM is a
compressed cohomology algorithm that processes inserted simplices through boundary
annotations, whereas the coreference experiment here is a different dual recurrence.

\subsection{Reproducibility is now part of the experiment}
The fair profile-selection validation is no longer just a list of shell commands in the
notebook history.  The fair-validation driver records the validation splits, timing
parameters, candidate gate, profile metrics, artifact paths, and command recipe in a
Markdown manifest.  This manifest should be regenerated whenever the validation CSVs are
regenerated.

\subsection{Limitations}

The present experiments are useful for choosing the next technical direction, but they
should not yet be read as final performance claims.
\begin{itemize}
  \item The Roadmap comparison uses selected truncated Rips objects, not a full
        reproduction of the Roadmap benchmark study.
  \item GUDHI is used as a strong external persistence baseline on the same filtered
        complexes.  It is a CAM-style comparison point, not a direct timing of the
        original CAM implementation.
  \item The native GUDHI-view result is encouraging as an integration route but remains
        below parity end-to-end in the Release benchmark.  The direct view is much faster
        than compact import, and the reducer-only ratios are favorable, but pre-reducer
        construction costs dominate.
  \item Several synthetic timings are in the microsecond range.  At that scale, small
        changes in machine state can change the measured winner, so selector conclusions
        must be based on aggregate behavior and hold-out validation.
  \item The current C++ core is optimized enough to compare strategies, but it is not yet
        an optimized large-complex construction library.
  \item The profile selector is still experimental.  The structured image-grid branch is
        now the leading candidate to validate, but the current evidence does not yet
        justify promoting it to the default API policy.
  \item Perseus export is kept as an optional barcode-validation path.  We do not use
        Perseus timings in the main comparison because command-line/file-I/O timings would
        not be directly comparable with the in-process benchmark timings.
\end{itemize}

\subsection{Discussion questions for Gilles Bertrand}

The current results suggest the following discussion questions.
\begin{enumerate}
  \item Is annotation volume a more appropriate objective than critical-cell count for the
        kind of Morse persistence pipeline tested here?
  \item Can one predict, from local combinatorics of a filtration level, whether a
        same-level sequence will keep reference annotations small?
  \item Which Morse sequence construction should be implemented next in the shared
        strategy interface: another theoretical Max/Min variant, a refined flooding rule,
        or a more global plateau method?
  \item What would constitute a fair direct comparison with CAM: a new implementation of
        CAM, a comparison with an existing cohomology implementation, or a benchmark framed
        only around compressed annotations?
  \item Are structured plateau examples, such as terrain and image-grid filtrations,
        central enough to drive the selector design, or should
        \texttt{adaptive-structured} remain an explicitly experimental option?
  \item What additional native GUDHI-view benchmarks would make the integration evidence
        convincing: more flag/clique examples, larger grid filtrations, Roadmap-style Rips
        objects constructed directly in a \texttt{Gudhi::Simplex\_tree}, or all of these?
  \item Which additional benchmark families would best reduce overfitting to the current
        random, Roadmap, terrain, and image-grid splits?
\end{enumerate}

\subsection{Current recommendation}

The recommended next technical step is to keep the original adaptive selector with the
family-aware gate as the experimental control, use the decision-summary artifact as the
selector audit, and promote \texttt{adaptive-structured} to the next validation candidate
before changing the default user-facing API.  In parallel, the native GUDHI-view benchmark
should be broadened because the current Release native run shows partial parity on
structured grids and also identifies concrete pre-reducer bottlenecks that still prevent a
uniform performance claim.
In particular, the experimental pipeline should:
\begin{enumerate}
  \item keep \texttt{family-aware} plus \texttt{adaptive} as the practical experimental
        selector for now, while reporting the effective metric used in every CSV row;
  \item validate \texttt{adaptive-structured} on fresh image-grid, terrain, plateau, and
        Rips splits, because the current aggregate audit removes the image-grid gap when
        that branch is treated as the active metric;
  \item use \texttt{run\_fair\_profile\_validation.py} as the canonical validation entry
        point, so the CSVs, summaries, generated LaTeX fragments, and manifest stay
        synchronized;
  \item continue expanding hold-out validation before replacing the current default
        selector, especially on structured grid/image-like filtrations where the current
        evidence favors a structured branch but with scale-dependent margins;
  \item keep the plateau-focused validation preset as a diagnostic control; the current
        run does not justify switching plateau cases to \texttt{profile-total-work};
  \item treat \texttt{terrain-holdout}, \texttt{image-holdout}, and
        \texttt{image-stress} as controls for any future structured branch before it
        becomes an API default;
  \item use the shape and plateau fields, including level concentration, entropy,
        largest-level size, and simplex-to-vertex ratio, to learn or tabulate selector
        rules, but treat tabulated gates as diagnostics until they survive fresh
        validation;
  \item keep Perseus export as an optional barcode-validation path rather than a timing
        baseline at this stage;
  \item keep GUDHI as the main external baseline for the current simplicial experiments,
        and extend the native \texttt{Simplex\_tree} validation beyond the current
        structured-grid and flag validation cases;
  \item keep the sequence-constructor interface open so future algorithms, including
        flooding variants and newer Morse-sequence schemes, can be added to the same
        profile-and-select loop.
\end{enumerate}

\begin{thebibliography}{9}

\bibitem{cam}
J.-D. Boissonnat, T. K. Dey, and C. Maria.
\newblock The Compressed Annotation Matrix: an Efficient Data Structure for Computing
Persistent Cohomology.
\newblock arXiv:1304.6813, 2013.

\bibitem{roadmap}
N. Otter, M. A. Porter, U. Tillmann, P. Grindrod, and H. A. Harrington.
\newblock A roadmap for the computation of persistent homology.
\newblock \emph{EPJ Data Science}, 6, 2017.

\end{thebibliography}

\end{document}
